\documentclass{article}

\newcounter{listing}

% ==============================================================================
% LOAD REFSUITE PACKAGE WITH EXPL3 KEY CONFIGURATION
% ==============================================================================
\usepackage{refsuite}
\SetReferencesConfiguration{
  hyperref = true,
  hyperref/options = { unicode, hyperindex, pdfpagelabels },
  hyperref/setup = {
    breaklinks,
    colorlinks,
    linktocpage,
    urlcolor = blue,
    linkcolor = red,
    filecolor = magenta,
    citecolor = green,
    pdftitle = {RefSuite expl3 Flow Demonstration},
    pdfauthor = {Jane Doe}
  },
  shortrefs = {
    picref = Fig.,
    tabref = Table,
    secref = Section,
    lstref = Listing
  },
  biblatex = true,
  biblatex/options = { backend = biber, style = numeric, defernumbers = true },
  biblatex/path = .,
  biblatex/bibliographies = {
    {
      name = main,
      title = References,
      files = example.refsuite.expl3-main.bib
    },
    {
      name = online,
      title = {Online Resources},
      files = example.refsuite.expl3-online.bib
    }
  },
  indexes = true,
  indexes/options = { makeindex, noautomatic },
  indexes/list = {
    { name = notation, title = {List of Notations}, intoc, columns = 2, columnsep = 20pt },
    { title = {Subject Index}, intoc, columns = 3, columnsep = 15pt }
  }
}

\begin{document}

\title{RefSuite expl3 Configuration Flow}
\author{Jane Doe}
\date{\today}
\maketitle

\tableofcontents

\section{Introduction}
\label{sec:introduction}

This example demonstrates the TeX-side expl3 key configuration flow for
\texttt{refsuite}. The configuration is written directly in the document preamble
with \verb|\SetReferencesConfiguration|.

The example uses the small bibliography files stored beside this source file, so
it can be compiled without external data.  It should be built from the
\texttt{examples} directory; generated auxiliary files are placed in
\texttt{.build/}, while the final PDF stays in this directory.

\subsection{Package Configuration Snippet}

The package and its supported features are configured through expl3 keys:

This snippet is intentionally close to the real preamble above.  The
\verb|biblatex/path = .| setting means bibliography resources are resolved from
the current \texttt{examples} directory, matching the local \texttt{latexmkrc}
output convention.

\begin{quote}
\begin{verbatim}
\usepackage{refsuite}
\SetReferencesConfiguration{
  hyperref = true,
  hyperref/setup = {
    colorlinks,
    linkcolor = red,
    citecolor = green
  },
  shortrefs = {
    picref = Fig.,
    tabref = Table,
    secref = Section,
    lstref = Listing
  },
  biblatex = true,
  biblatex/options = { backend = biber, style = numeric, defernumbers = true },
  biblatex/path = .,
  biblatex/bibliographies = {
    { name = main, title = References, files = example.refsuite.expl3-main.bib },
    { name = online, title = {Online Resources}, files = example.refsuite.expl3-online.bib }
  },
  indexes = true,
  indexes/options = { makeindex, noautomatic },
  indexes/list = {
    { name = notation, title = {List of Notations}, intoc },
    { title = {Subject Index}, intoc }
  }
}
\end{verbatim}
\end{quote}
\refstepcounter{listing}\label{lst:configuration}

\section{Main Content}

\subsection{Configured Short References}

The expl3 flow creates the same short-reference commands: \picref{fig:sample},
\tabref{tab:sample}, \secref{sec:conclusion}, and \lstref{lst:configuration}.
The lowercase commands use the configured text exactly, while the capitalized
commands are generated automatically for sentence starts.

\begin{figure}[htbp]
  \centering
  \fbox{\parbox{0.5\textwidth}{\centering Example figure}}
  \caption{Example figure with caption}
  \label{fig:sample}
\end{figure}

\begin{table}[htbp]
  \centering
  \begin{tabular}{|c|c|}
    \hline
    Column 1 & Column 2 \\
    \hline
    Data 1 & Data 2 \\
    Data 3 & Data 4 \\
    \hline
  \end{tabular}
  \caption{Example table}
  \label{tab:sample}
\end{table}

\subsection{Bibliography}

The expl3 key flow can configure multiple named bibliography groups. This
example cites books \cite{knuth1984texbook,lamport1994latex} and an online
resource \cite{ctan2024latex}.

Because two bibliography groups are configured, \texttt{refsuite} generates
separate printing commands.  They are used near the end of the document as
\verb|\printbibliographymain| and \verb|\printbibliographyonline|.

\subsection{Indexes}

The index configuration creates one named notation index and one ordinary subject
index.  Named indexes also receive generated print commands, so the notation
index can be printed with \verb|\printindexnotation|.

\index[notation]{$\gamma$} Gamma notation
\index[notation]{$\delta$} Delta notation
\index{LaTeX} Typesetting system
\index{bibliography} List of references
\index{expl3 configuration} expl3 configuration

\section{Conclusion}
\label{sec:conclusion}

The expl3 flow is useful when the package configuration should stay in the TeX
preamble and use the same key-value syntax as the package internals.

\printbibliographymain
\printbibliographyonline

\printindexnotation
\printindex

\end{document}
